\documentclass[11pt]{article}
\usepackage[margin=1in]{geometry}
\usepackage{amsmath, amssymb, amsthm, bm, mathtools}
\usepackage{booktabs, array}
\usepackage{enumitem}
\usepackage{siunitx}
\usepackage{hyperref}
\hypersetup{colorlinks=true, linkcolor=blue, urlcolor=blue, citecolor=blue}
\sisetup{round-mode=places,round-precision=3,group-separator={,}}

\newcommand{\dd}{\,\mathrm{d}}
\newcommand{\R}{\mathbb{R}}


\title{Solutions -- Practice Problem Set 1}
\author{ECON 308 -- International Economic Relations}
\date{Fall 2025}

\begin{document}
\maketitle



\begin{enumerate}[leftmargin=*]
\item \textbf{Gravity intuition.} Larger GDPs imply more production and income, scaling the number of potential buyers and sellers; microfoundations include love-of-variety models and market-size effects that raise the mass of matches. Distance proxies trade costs (freight, time-to-market, information frictions), which depress the intensive and extensive margins of trade. Search/information models also imply distance lowers match probabilities, while per-unit transport costs directly reduce delivered profitability.





\end{enumerate}




\begin{enumerate}[leftmargin=*,resume]

\item \textbf{Relative wages under free trade.} With perfect competition and specialization, each country's wage equals the value of marginal product in the sector it produces: $w=\dfrac{P_g}{a_{Lg}}$. The \emph{relative} wage across countries is pinned down by world prices and comparative advantage: the high-wage country must have sufficiently high productivity in its export good to remain competitive at world prices.


\item \textbf{Unit labor costs and prices.} In autarky, competitive pricing implies $P_g=w\,a_{Lg}$, so relative prices equal unit labor cost ratios. With trade, world prices need not equal a country's own unit costs; instead, specialization ensures the country produces the good where it has the lowest unit cost at world prices, and wages adjust to equalize unit costs in the produced sector.
\end{enumerate}




\begin{enumerate}[leftmargin=*, resume]
\item \textbf{Gravity Model: Baseline Predictions and Elasticities.}

\emph{Given:} $T_{ij}=10^{-6}\dfrac{Y_iY_j}{D_{ij}}$, with $(Y_A,Y_B,Y_C)=(2000,1000,500)$ (billions), and distances in km:
$D_{AB}=6{,}000$, $D_{AC}=1{,}500$, $D_{BC}=2{,}000$.

\begin{enumerate}
\item \emph{Predicted bilateral exports.}
\[
\begin{aligned}
T_{AB}&=10^{-6}\frac{(2000)(1000)}{6000}
=10^{-6}\cdot \frac{2{,}000{,}000}{6000}
=10^{-6}\cdot 333.\overline{3}
=0.000333\ \text{(billions)}\\
&\Rightarrow \ \ \$0.333\ \text{million.}\\[6pt]
T_{AC}&=10^{-6}\frac{(2000)(500)}{1500}
=10^{-6}\cdot \frac{1{,}000{,}000}{1500}
=10^{-6}\cdot 666.\overline{6}
=0.000667\ \text{(billions)}\\
&\Rightarrow \ \ \$0.667\ \text{million.}\\[6pt]
T_{BC}&=10^{-6}\frac{(1000)(500)}{2000}
=10^{-6}\cdot \frac{500{,}000}{2000}
=10^{-6}\cdot 250
=0.000250\ \text{(billions)}\\
&\Rightarrow \ \ \$0.250\ \text{million.}
\end{aligned}
\]

\item \emph{Effect of a 10\% fall in distance on $T_{AB}$.} The log-specification gives the elasticity of trade with respect to distance as $-\gamma=-1$. Exactly:
\[
\frac{T_{AB}^{\,\text{new}}}{T_{AB}^{\,\text{old}}}
=\left(\frac{D_{AB}^{\text{new}}}{D_{AB}^{\text{old}}}\right)^{-\gamma}
=(0.9)^{-1}=\frac{1}{0.9}=1.\overline{1}
\]
So $T_{AB}$ rises by $11.\overline{1}\%$.

\item \emph{Effect of 5\% GDP increases in both A and B.} Because $T_{ij}\propto Y_iY_j$,
\[
\frac{T_{AB}^{\,\text{new}}}{T_{AB}^{\,\text{old}}}=(1.05)(1.05)=1.1025 \Rightarrow \text{increase of } 10.25\%.
\]

\item \emph{Log residual against observed \$1.3 bn.} 
\[
\ln T^{\text{obs}}_{AB}-\ln T^{\text{pred}}_{AB}
=\ln(1.3)-\ln(0.000333\ldots)\approx 8.269.
\]
A positive residual indicates $AB$ trade exceeds gravity prediction, consistent with omitted frictions/facilitators (e.g., common language, trade agreement, high-quality transport links, or substantial services trade not captured by simple distance).
\end{enumerate}




\end{enumerate}





\begin{enumerate}[leftmargin=*, resume]
\item \textbf{Two-Country, Two-Good Ricardian Basics.}

\emph{Given:} Unit labor requirements
\[
\begin{array}{@{}lcc@{}}
\toprule
 & a_{LW} & a_{LC}\\
\midrule
H & 2 & 4\\
F & 6 & 3\\
\bottomrule
\end{array}
\qquad L_H=L_F=240.
\]

\begin{enumerate}
\item \emph{Opportunity costs and comparative advantage.} In the Ricardian model,
\[
\text{OC}_H(W \text{ in } C)=\frac{a_{LW}}{a_{LC}}=\frac{2}{4}=0.5,\quad
\text{OC}_H(C \text{ in } W)=\frac{a_{LC}}{a_{LW}}=\frac{4}{2}=2.
\]
\[
\text{OC}_F(W \text{ in } C)=\frac{6}{3}=2,\quad
\text{OC}_F(C \text{ in } W)=\frac{3}{6}=0.5.
\]
Home (H) has lower OC in $W$ ($0.5<2$), so $H$ has comparative advantage in $W$; Foreign (F) in $C$.

\item \emph{PPFs and autarky prices.} PPFs satisfy $a_{LW}W+a_{LC}C\le L$.
\[
\text{H: } 2W+4C\le 240 \Rightarrow W=120-2C\ (\text{slope }-2).
\qquad
\text{F: } 6W+3C\le 240 \Rightarrow W=40-\tfrac{1}{2}C\ (\text{slope }-\tfrac{1}{2}).
\]
Under perfect competition, autarky relative price equals the unit-cost ratio:
\[
\left(\frac{P_W}{P_C}\right)^{A}_H=\frac{a_{LW}}{a_{LC}}=\tfrac{1}{2}=0.5,
\qquad
\left(\frac{P_W}{P_C}\right)^{A}_F=\frac{6}{3}=2.
\]

\item \emph{Range for complete specialization.} With $H$ in $W$ and $F$ in $C$,
\[
0.5 < \frac{P_W}{P_C} < 2.
\]

\item \emph{Specialization at $P_W/P_C=1.2$.} Since $0.5<1.2<2$, $H$ specializes in $W$, $F$ in $C$.
\[
\text{H output: } W_H=\frac{L_H}{a_{LW}}=\frac{240}{2}=120,\quad C_H=0.
\qquad
\text{F output: } C_F=\frac{L_F}{a_{LC}}=\frac{240}{3}=80,\quad W_F=0.
\]
$H$ exports $W$ and imports $C$; $F$ the reverse.

\item \emph{Real wages under free trade.} With specialization, $w_H=\dfrac{P_W}{a_{LW}}$ and $w_F=\dfrac{P_C}{a_{LC}}$.
\[
\text{H:}\quad \frac{w_H}{P_W}=\frac{1}{a_{LW}}=\frac{1}{2}=0.5\ \text{(units of $W$)},\quad
\frac{w_H}{P_C}=\frac{P_W/P_C}{a_{LW}}=\frac{1.2}{2}=0.6\ \text{(units of $C$)}.
\]
\[
\text{F:}\quad \frac{w_F}{P_C}=\frac{1}{a_{LC}}=\frac{1}{3}\ \text{(units of $C$)},\quad
\frac{w_F}{P_W}=\frac{1}{a_{LC}}\cdot\frac{P_C}{P_W}=\frac{1}{3\cdot 1.2}=\frac{1}{3.6}\approx 0.278\ \text{(units of $W$)}.
\]
\end{enumerate}



\item \textbf{Terms of Trade and Real Incomes (price falls to $P_W/P_C=0.8$).}
\begin{enumerate}
\item Since $0.5<0.8<2$, specialization pattern is unchanged: $H\to W$, $F\to C$.

\item \emph{Real wages at $P_W/P_C=0.8$.}
\[
\text{H:}\ \frac{w_H}{P_W}=\frac{1}{2}=0.5\ \text{W},\quad
\frac{w_H}{P_C}=\frac{0.8}{2}=0.4\ \text{C}.
\quad
\text{F:}\ \frac{w_F}{P_C}=\frac{1}{3}\ \text{C},\quad
\frac{w_F}{P_W}=\frac{1}{3\cdot 0.8}=\frac{1}{2.4}\approx 0.417\ \text{W}.
\]
\emph{Comparison to autarky:} Autarky real wages are $H:(1/2\ \text{W},\ 1/4\ \text{C})$; $F:(1/6\ \text{W},\ 1/3\ \text{C})$. Both countries gain from trade relative to autarky.

\item \emph{Terms of trade (TOT) effects.} $H$ exports $W$, so its TOT is $P_W/P_C$. When this falls from $1.2$ to $0.8$, $H$'s purchasing power in $C$ falls ($0.6\to 0.4$), while $F$'s TOT $=(P_C/P_W)$ improves (its real wage in $W$ rises $0.278\to 0.417$).
\end{enumerate}

\item \textbf{World Relative Supply with Three Countries.}

\emph{Given:}
\[
\begin{array}{@{}lcc@{}}
\toprule
 & a_{LW} & a_{LC}\\
\midrule
1 & 1 & 2\\
2 & 2 & 3\\
3 & 4 & 5\\
\bottomrule
\end{array}
\qquad L_1=L_2=L_3=100.
\]
Compute $a_{LW}/a_{LC}$: Country 1: $0.5$; 2: $2/3\approx 0.667$; 3: $4/5=0.8$.

\begin{enumerate}
\item \emph{Ordering by comparative advantage in $W$:} $1$ (best), $2$, $3$.

\item \emph{World relative supply (WRS).} Each country produces $W$ if $P_W/P_C\ge a_{LW}/a_{LC}$; otherwise $C$.
\[
\begin{array}{l|c|c|c}
\text{Price range} & \text{W producers} & \text{W supply} & \text{C supply}\\
\hline
P<0.5 & \text{none} & 0 & \frac{100}{2}+\frac{100}{3}+\frac{100}{5}=50+33.\overline{3}+20=103.\overline{3}\\
0.5<P<2/3 & 1 & \frac{100}{1}=100 & \frac{100}{3}+\frac{100}{5}=33.\overline{3}+20=53.\overline{3}\\
2/3<P<0.8 & 1,2 & \frac{100}{1}+\frac{100}{2}=150 & \frac{100}{5}=20\\
P>0.8 & 1,2,3 & \frac{100}{1}+\frac{100}{2}+\frac{100}{4}=175 & 0
\end{array}
\]
Thus the step values of $RS^{W/C}\equiv W/C$ are: $0$; $\dfrac{100}{53.\overline{3}}\approx 1.875$; $\dfrac{150}{20}=7.5$; and $+\infty$.
Kinks at $P=0.5,\ 2/3,\ 0.8$.

\item \emph{Equilibrium with $RD^{W/C}=1.1-0.2(P_W/P_C)$.}
Because $RD$ is continuous and decreasing, and $RS$ is a step function, the market clears either (i) on a step (if $RD$ hits a step value), or (ii) at a kink with the marginal country partially specialized. Evaluate $RD$ at kinks:
\[
RD(0.5)=1.1-0.2(0.5)=1.0,\quad RD(2/3)\approx 0.967,\quad RD(0.8)=0.94.
\]
All are \emph{below} the adjacent step value $1.875$. Hence equilibrium occurs at the \emph{lowest} kink $P_W/P_C=0.5$, with Country 1 splitting labor between $W$ and $C$ to make $RS=RD=1$.

Let $x\in[0,1]$ be the share of Country 1's labor in $W$ at $P=0.5$.
\[
W=\frac{100x}{1}=100x,
\qquad
C=\frac{100(1-x)}{2}+\frac{100}{3}+\frac{100}{5}=50(1-x)+53.\overline{3}=103.\overline{3}-50x.
\]
Set $W/C=1$:
\[
100x=103.\overline{3}-50x \ \Rightarrow\ 150x=103.\overline{3}\ \Rightarrow\ x\approx 0.6889.
\]
\emph{Production at equilibrium:} Country 1 uses $\approx 68.9\%$ of $L_1$ in $W$; Countries 2 and 3 produce only $C$.
\end{enumerate}

\end{enumerate}




\end{document}
